\documentclass[11pt,letterpaper]{article}
\usepackage{cornell-notes}
\usepackage{needspace}

\title{Chapter 2: Design Review}
\author{}
\date{}

\setDocTitle{Chapter 2: Design Review}
\setDocSubtitle{Application Security Review}
\setCornellCollection{Software Security Assessment}
\setCornellUnitType{Chapter}
\setCornellUnitNumber{2}
\setCornellUnitTitle{Design Review}
\setCornellCompanionLabel{Study and audit reference}
\setCornellSourceRange{pp. 25--66}
\begin{document}
\maketitle

\begin{center}
\small\color{Meta}\textbf{Coverage range:} pp. 25--66 \quad\textbullet\quad \textbf{Format:} Cornell cue-and-notes
\end{center}
\vspace{0.25em}
\begin{CornellOverviewBox}{Review Orientation}
\textbf{Central problem.} Evaluate whether a software design can enforce its intended security policy before implementation details distract from structural weaknesses.
\textbf{Why it matters.} Security reviewers use this material to connect attacker-controlled conditions to violated invariants, consequential operations, and defensible remediation.
\textbf{Prerequisites.} Security policies, trust boundaries, basic software architecture, and the confidentiality-integrity-availability model.
\textbf{Assessment relationship.} It turns the foundational vocabulary into a structured architecture and threat-modeling practice that precedes implementation review.
\end{CornellOverviewBox}
\section{Learning Objectives}
\begin{itemize}
\item Explain the security significance of algorithms.
\item Identify attacker influence and failed assumptions associated with abstraction and decomposition.
\item Trace security-relevant data or state through trust relationships.
\item Evaluate whether controls addressing least privilege and separation preserve the intended invariant.
\item Audit implementation and error paths related to complete mediation.
\item Distinguish safe and unsafe uses of fail-safe defaults and economy.
\item Diagnose a failure involving authentication from concrete evidence.
\item Verify remediation and test coverage for authorization.
\end{itemize}
\section{Cornell Cue-and-Notes}
\Needspace{8\baselineskip}
\subsection{Software Design Fundamentals}
\begin{longtable}{>{\raggedright\arraybackslash\columncolor{CornellCueBackground}}p{0.285\textwidth}|>{\raggedright\arraybackslash}p{0.625\textwidth}}
\rowcolor{Navy}\color{white}\textbf{Cue / Recall Prompt} & \color{white}\textbf{Technical Notes and Audit Reasoning} \\
\endfirsthead
\rowcolor{Navy}\color{white}\textbf{Cue / Recall Prompt} & \color{white}\textbf{Technical Notes and Audit Reasoning} \\
\endhead
\term{Algorithms}\\[-0.1em]{\small Can a correct algorithm still be insecure?} & An algorithm may produce functionally correct output while exposing secrets, accepting ambiguous input, exhausting resources, or relying on assumptions an attacker can violate. Security review examines adversarial behavior and boundary conditions in addition to nominal correctness.\par\smallskip\textbf{Audit move:} Test algorithm assumptions against malicious sizes, ordering, repetition, and state transitions. \\\hline
\term{Abstraction and decomposition}\\[-0.1em]{\small How can component boundaries hide security flaws?} & Decomposition makes systems understandable but can distribute one security decision across several components. Each abstraction may be locally correct while the composition leaves a gap, duplicates enforcement inconsistently, or loses context.\par\smallskip\textbf{Audit move:} Trace each security decision end to end across every participating component. \\\hline
\term{Trust relationships}\\[-0.1em]{\small What should justify a trust edge?} & Trust must follow from an enforceable property such as authenticated identity, protected execution, or validated data--not from location, naming, or expected behavior alone. Transitive trust deserves special scrutiny because compromise propagates.\par\smallskip\textbf{Audit move:} Document why each component or principal is trusted and what happens if it is malicious. \\\hline
\end{longtable}
\begin{CornellSummaryBox}{Section Summary}
Algorithms, decomposition, and trust relationships determine which guarantees an implementation can realistically preserve.
\end{CornellSummaryBox}
\Needspace{8\baselineskip}
\subsection{Principles and Fundamental Flaws}
\begin{longtable}{>{\raggedright\arraybackslash\columncolor{CornellCueBackground}}p{0.285\textwidth}|>{\raggedright\arraybackslash}p{0.625\textwidth}}
\rowcolor{Navy}\color{white}\textbf{Cue / Recall Prompt} & \color{white}\textbf{Technical Notes and Audit Reasoning} \\
\endfirsthead
\rowcolor{Navy}\color{white}\textbf{Cue / Recall Prompt} & \color{white}\textbf{Technical Notes and Audit Reasoning} \\
\endhead
\term{Least privilege and separation}\\[-0.1em]{\small How should privilege be distributed?} & Give each principal only the authority needed for its task and separate incompatible powers. Narrow privileges reduce blast radius, while separation of duties prevents one compromised element from completing a sensitive action alone.\par\smallskip\textbf{Audit move:} Compare every privilege grant with the exact operation and lifetime that require it. \\\hline
\term{Complete mediation}\\[-0.1em]{\small Where must policy checks occur?} & Every security-sensitive access must be authorized using current, trustworthy state. Caching or performing checks only at entry points can create gaps when objects, identities, or policy change before use.\par\smallskip\textbf{Audit move:} Locate all access paths to protected resources and verify a non-bypassable decision point. \\\hline
\term{Fail-safe defaults and economy}\\[-0.1em]{\small What happens when information is missing or a component fails?} & Access should be denied unless explicitly allowed, and the enforcement design should remain small enough to understand and test. Complex exception logic and permissive fallback behavior enlarge the attack surface.\par\smallskip\textbf{Audit move:} Inspect initialization, recovery, timeout, and error paths for unintended grants. \\\hline
\end{longtable}
\begin{CornellSummaryBox}{Section Summary}
Secure designs minimize authority, centralize policy, make defaults safe, and continue to protect assets when components fail.
\end{CornellSummaryBox}
\Needspace{8\baselineskip}
\subsection{Enforcing Security Policy}
\begin{longtable}{>{\raggedright\arraybackslash\columncolor{CornellCueBackground}}p{0.285\textwidth}|>{\raggedright\arraybackslash}p{0.625\textwidth}}
\rowcolor{Navy}\color{white}\textbf{Cue / Recall Prompt} & \color{white}\textbf{Technical Notes and Audit Reasoning} \\
\endfirsthead
\rowcolor{Navy}\color{white}\textbf{Cue / Recall Prompt} & \color{white}\textbf{Technical Notes and Audit Reasoning} \\
\endhead
\term{Authentication}\\[-0.1em]{\small What evidence binds an action to an identity?} & Authentication establishes a claimed identity with evidence appropriate to the threat. The design must protect credentials, enrollment, recovery, session continuity, and transitions between anonymous and authenticated state.\par\smallskip\textbf{Audit move:} Model credential lifecycle, recovery, reauthentication, and session-binding paths. \\\hline
\term{Authorization}\\[-0.1em]{\small Is an authenticated principal allowed to perform this action?} & Authorization evaluates subject, object, action, and context against policy. Confusing authentication with authorization or placing decisions only in presentation layers produces horizontal and vertical access-control failures.\par\smallskip\textbf{Audit move:} Build a subject-object-action matrix and map every decision to an enforcement point. \\\hline
\term{Accountability}\\[-0.1em]{\small Can sensitive actions be reconstructed reliably?} & Accountability requires trustworthy records with identity, action, target, result, and time. Logs must resist tampering and avoid becoming a new source of secret disclosure.\par\smallskip\textbf{Audit move:} Verify coverage, integrity, retention, access control, and failure behavior of audit records. \\\hline
\term{Confidentiality, integrity, availability}\\[-0.1em]{\small How do architectural controls preserve core properties?} & Confidentiality needs controlled disclosure, integrity needs authorized and attributable modification, and availability needs resilience against failure and resource abuse. Tradeoffs should be explicit because a control for one property can weaken another.\par\smallskip\textbf{Audit move:} For each critical asset, record required properties, threats, controls, and residual risk. \\\hline
\end{longtable}
\begin{CornellSummaryBox}{Section Summary}
Authentication, authorization, accountability, confidentiality, integrity, and availability must be designed as coordinated controls rather than isolated features.
\end{CornellSummaryBox}
\Needspace{8\baselineskip}
\subsection{Threat Modeling}
\begin{longtable}{>{\raggedright\arraybackslash\columncolor{CornellCueBackground}}p{0.285\textwidth}|>{\raggedright\arraybackslash}p{0.625\textwidth}}
\rowcolor{Navy}\color{white}\textbf{Cue / Recall Prompt} & \color{white}\textbf{Technical Notes and Audit Reasoning} \\
\endfirsthead
\rowcolor{Navy}\color{white}\textbf{Cue / Recall Prompt} & \color{white}\textbf{Technical Notes and Audit Reasoning} \\
\endhead
\term{Information collection}\\[-0.1em]{\small What must be known before identifying threats?} & Reviewers need component responsibilities, data classifications, identities, privileges, trust boundaries, entry points, protocols, persistence, deployment topology, and operational dependencies. Missing architecture information is itself a review risk.\par\smallskip\textbf{Audit move:} Obtain or reconstruct data-flow, deployment, and privilege views before ranking threats. \\\hline
\term{Application architecture modeling}\\[-0.1em]{\small Which model makes security relationships visible?} & Useful models show processes, data stores, flows, external actors, boundaries, and policy-enforcement points. The goal is analytical accuracy, not decorative completeness.\par\smallskip\textbf{Audit move:} Validate the model against code, configuration, and runtime observations. \\\hline
\term{Threat identification and documentation}\\[-0.1em]{\small How are threats made actionable?} & For each asset and boundary, describe attacker capability, precondition, action, violated assumption, impact, and candidate controls. Findings should distinguish evidence from uncertainty and identify the implementation areas that need deeper review.\par\smallskip\textbf{Audit move:} Write testable threat statements and link each to owners, evidence needs, and review targets. \\\hline
\term{Prioritization}\\[-0.1em]{\small Which implementation areas should be reviewed first?} & Prioritize code that combines attacker reachability, high privilege, sensitive assets, complex parsing, dangerous operations, or weak compensating controls. Ranking should reflect plausible impact and uncertainty, not just feature size.\par\smallskip\textbf{Audit move:} Use the threat model to create a risk-ordered implementation-review queue. \\\hline
\end{longtable}
\begin{CornellSummaryBox}{Section Summary}
Threat modeling converts architecture knowledge into prioritized attacker stories and targeted implementation-review questions.
\end{CornellSummaryBox}
\begin{CornellLegacyBox}{Historical Context}
Specific authentication and cryptographic mechanisms evolve; retain the chapter's design principles while evaluating concrete mechanisms against current platform guidance during a real assessment.
\end{CornellLegacyBox}
\section{Security-Review Checklist}
\begin{CornellChecklist}
\item \textbf{Scoping}
Test algorithm assumptions against malicious sizes, ordering, repetition, and state transitions.
\item \textbf{Attack-surface identification}
Trace each security decision end to end across every participating component.
\item \textbf{Input and data-flow analysis}
Document why each component or principal is trusted and what happens if it is malicious.
\item \textbf{Trust-boundary review}
Compare every privilege grant with the exact operation and lifetime that require it.
\item \textbf{Candidate-point inspection}
Locate all access paths to protected resources and verify a non-bypassable decision point.
\item \textbf{Control-flow review}
Inspect initialization, recovery, timeout, and error paths for unintended grants.
\item \textbf{Resource and lifetime review}
Model credential lifecycle, recovery, reauthentication, and session-binding paths.
\item \textbf{Error-path analysis}
Build a subject-object-action matrix and map every decision to an enforcement point.
\item \textbf{Mitigation validation}
Verify coverage, integrity, retention, access control, and failure behavior of audit records.
\item \textbf{Evidence and reporting}
For each critical asset, record required properties, threats, controls, and residual risk.
\item \textbf{Scoping}
Obtain or reconstruct data-flow, deployment, and privilege views before ranking threats.
\item \textbf{Attack-surface identification}
Validate the model against code, configuration, and runtime observations.
\end{CornellChecklist}
\section{Chapter Synthesis}
Design review asks whether the system can remain secure even when inputs are malicious and components fail. The reviewer models architecture, makes trust explicit, verifies that policy properties have complete enforcement paths, and converts threats into prioritized code-review targets. Fundamental flaws deserve priority because local patches rarely repair a broken trust or authorization model.
\section{Self-Test Questions}
\begin{enumerate}
\item Can a correct algorithm still be insecure?
\item How can component boundaries hide security flaws?
\item What should justify a trust edge?
\item How should privilege be distributed?
\item Where must policy checks occur?
\item What happens when information is missing or a component fails?
\item \textbf{Scenario:} A user interface hides an administrator action, but the service endpoint accepts the request from ordinary users. Which design principle failed?
\item \textbf{Scenario:} A service caches an authorization result and continues using it after account privileges are revoked. What should the design review examine?
\end{enumerate}
\clearpage
\subsection*{Answer Key}
\begin{enumerate}
\item An algorithm may produce functionally correct output while exposing secrets, accepting ambiguous input, exhausting resources, or relying on assumptions an attacker can violate. Security review examines adversarial behavior and boundary conditions in addition to nominal correctness.
\item Decomposition makes systems understandable but can distribute one security decision across several components. Each abstraction may be locally correct while the composition leaves a gap, duplicates enforcement inconsistently, or loses context.
\item Trust must follow from an enforceable property such as authenticated identity, protected execution, or validated data--not from location, naming, or expected behavior alone. Transitive trust deserves special scrutiny because compromise propagates.
\item Give each principal only the authority needed for its task and separate incompatible powers. Narrow privileges reduce blast radius, while separation of duties prevents one compromised element from completing a sensitive action alone.
\item Every security-sensitive access must be authorized using current, trustworthy state. Caching or performing checks only at entry points can create gaps when objects, identities, or policy change before use.
\item Access should be denied unless explicitly allowed, and the enforcement design should remain small enough to understand and test. Complex exception logic and permissive fallback behavior enlarge the attack surface.
\item Authorization was placed in a presentation layer instead of a complete, non-bypassable mediation point. The service must enforce subject-object-action policy independently of interface visibility.
\item Examine complete mediation, cache lifetime, revocation propagation, time-of-check/time-of-use gaps, and fail-safe behavior when current policy cannot be retrieved.
\end{enumerate}
\section{Key-Term Recap}
\begin{longtable}{>{\raggedright\arraybackslash}p{0.27\textwidth}p{0.65\textwidth}}
\toprule
\textbf{Term} & \textbf{Working meaning for review} \\
\midrule
\endfirsthead
\toprule
\textbf{Term} & \textbf{Working meaning for review} \\
\midrule
\endhead
Algorithms & An algorithm may produce functionally correct output while exposing secrets, accepting ambiguous input, exhausting resources, or relying on assumptions an attacker can violate. \\
Abstraction and decomposition & Decomposition makes systems understandable but can distribute one security decision across several components. \\
Trust relationships & Trust must follow from an enforceable property such as authenticated identity, protected execution, or validated data--not from location, naming, or expected behavior alone. \\
Least privilege and separation & Give each principal only the authority needed for its task and separate incompatible powers. \\
Complete mediation & Every security-sensitive access must be authorized using current, trustworthy state. \\
Fail-safe defaults and economy & Access should be denied unless explicitly allowed, and the enforcement design should remain small enough to understand and test. \\
Authentication & Authentication establishes a claimed identity with evidence appropriate to the threat. \\
Authorization & Authorization evaluates subject, object, action, and context against policy. \\
Accountability & Accountability requires trustworthy records with identity, action, target, result, and time. \\
Confidentiality, integrity, availability & Confidentiality needs controlled disclosure, integrity needs authorized and attributable modification, and availability needs resilience against failure and resource abuse. \\
\bottomrule
\end{longtable}
\section{One-Sentence Takeaway}
\begin{CornellExamBox}{One-Sentence Takeaway}
\textbf{A secure implementation cannot compensate for an architecture whose trust relationships and policy-enforcement paths are incomplete or internally inconsistent.}
\end{CornellExamBox}
\end{document}
